\chapter{Bounded curvature at bounded distance}
\label{chap:bounded-curvature-bounded-distance}
\section{Pinching toward positive: the definitions}\label{10.1}

\begin{definition}[Curvature pinched toward positive]
\label{def:curvature-pinched-toward-positive}
\uses{def:generalized-ricci-flow, def:curvature-operator}
\label{pinchdef}

Let $({\mathcal M},G)$ be a generalized three-dimensional Ricci flow
whose domain of definition is contained in $[0,\infty)$. For each
$x\in {\mathcal M}$, let $\nu(x)$ be the smallest eigenvalue of
$\Rm(x)$ on $\wedge^2T_xM_{{\bf t}(x)}$, as measured with
respect to a $G(x)$-orthonormal basis for the horizontal space at
$x$, and set $X(x)=\max(0,-\nu(x))$. We say that $({\mathcal
M},G)$ has curvature {\em pinched toward positive} if, for all $x\in
{\mathcal M}$, if the following two inequalities hold:
\begin{enumerate}
\item[(1)] $$R(x)\ge \frac{-6}{1+4{\bf t}(x)},$$
\item[(2)] $$R(x)\geq 2X(x)\left(\log X(x)+\log
(1+{\bf t}(x))-3\right),
$$ whenever
$0<X(x)$.
\end{enumerate}
\end{definition}

Theorem~\ref{pincha} gives preservation of these inequalities for
compact three-dimensional Ricci flows. In particular, if
$|\Rm(p,0)|\le 1$ initially, the flow is pinched toward positive.
Fix $\epsilon_0>0$ sufficiently small such that for any
$0<\epsilon\le \epsilon_0$ all the results of the Appendix hold for
$2\epsilon$ and $\alpha=10^{-2}$, and Proposition~\ref{narrows}
holds for $2\epsilon$.


\section{The statement of the theorem}

\begin{theorem}[Bounded curvature at bounded distance]
\label{thm:bounded-curvature-bounded-distance}
\uses{def:curvature-pinched-toward-positive, def:strong-epsilon-neck-canonical-neighborhood, nocones}
\label{bcbd}

Fix $0<\epsilon\le \epsilon_0$ and  $C<\infty$. Then for each
$A<\infty$ there are $D_0<\infty$ and $D<\infty$ depending on $A$,
$\epsilon$ and $C$ such that the following holds. Suppose that
$({\mathcal M},G)$ is a generalized three-dimensional Ricci flow
whose interval of definition is contained in $[0,\infty)$, and
suppose that $x\in {\mathcal M}$. Set $t={\bf t}(x)$. We suppose
that these data satisfy the following:
\begin{enumerate}
\item[(1)] $({\mathcal M},G)$ has curvature pinched toward positive.
\item[(2)] Every point $y\in{\mathcal  M}$ with
$R(y)\ge 4R(x)$  and ${\bf t}(y)\le t$ has a strong
$(C,\epsilon)$-canonical neighborhood.
\end{enumerate}
If $R(x)\ge D_0$, then $R(y)\le DR(x)$ for all $y\in
B(x,t,AR(x)^{-1/2})$.
\end{theorem}


This chapter is devoted to the proof of this theorem. The proof is
by contradiction. Suppose that there is some $A_0<\infty$ for which
the result fails. Then there are a sequence of generalized
three-dimensional Ricci flows $({\mathcal M}_n,G_n)$  whose
intervals of definition are contained in $[0,\infty)$ and whose
curvatures are pinched toward positive. Also, there are points
$x_n\in {\mathcal M}_n$ satisfying the second condition given in the
theorem and points $y_n\in {\mathcal M}_n$ such that for all $n$ we
have:
\begin{enumerate}
\item[(1)] $\lim_{n\rightarrow\infty}R(x_n)=\infty$.
\item[(2)] ${\bf t}(y_n)={\bf t}(x_n)$,
\item[(3)] $d(x_n,y_n)< A_0R(x_n)^{-1/2}$,
\item[(4)] $$\lim_{n\rightarrow\infty}\frac{R(y_n)}{R(x_n)}=\infty.$$
\end{enumerate}
The sequence above is fixed for the contradiction argument. Its
rescaled limit is an incomplete finite-diameter manifold
$U_\infty\cong S^2\times(0,1)$ with curvature bounded at one end and
unbounded at the other end ${\mathcal E}$; points sufficiently near
${\mathcal E}$ are centers of evolving $2\epsilon$-necks. Rescaling
near ${\mathcal E}$ gives both a cone limit and a non-flat smooth
Ricci-flow limit. Their scale ratio is bounded above and below, so
the latter is an open subset of a non-flat cone, contradicting
Theorem~\ref{nocones}.

\section{The incomplete geometric limit}\label{sectincomp}

Fix $({\mathcal M}_n,G_n,x_n)$ as above and shift and rescale it to
form the incomplete geometric limit.

We shift the time parameter of $({\mathcal M}_n,G_n)$ by $-{\bf
t}(x_n)$. We change notation and denote these shifted flows by
$({\mathcal M}_n,G_n)$. This allows us to arrange that ${\bf
t}(x_n)=0$ for all $n$. Since shifting leaves the curvature
unchanged, the shifted flows  satisfy a weaker version of curvature
pinched toward positive. Namely, for the shifted flows we have
\begin{eqnarray}\label{weakerpos}
R(x) & \ge & -6 \nonumber \\
R(x) & \ge & 2X(x)\left(\log(X(x))-3\right).
\end{eqnarray}

 We set $Q_n=R(x_n)$, and we denote by $M_n$ the $0$-time-slice of
${\mathcal M}_n$. We rescale $({\mathcal M}_n,G_n)$ by $Q_n$. Denote
by $({\mathcal M}'_n,G_n')$ the rescaled (and shifted) generalized
flows. For the rest of this argument we implicity use the metrics
$G'_n$. If we are referring to $G_n$ we mention it explicitly.

\subsection{The sequence of tubes}

 Let
$\gamma_n$ be a smooth path in $B_{G_n}(x_n,0,A_0Q_n^{-1/2})$ from
$x_n$ to $y_n$. For all $n$ sufficiently large we have
$R_{G'_n}(y_n) \gg 1$. Thus, there is a point $z_n\in \gamma_n$ such
that $R_{G'_n}(z_n)=4$ and such that on the sub-path
$\gamma_n|_{[z_n,y_n]}$ we have $R_{G'_n}\ge 4$. We replace
$\gamma_n$ by this sub-path. Now, with this replacement, according
to the second condition in the statement of the theorem, every point
of $\gamma_n$ has a strong $(C,\epsilon)$ canonical
neighborhood. As $n$ tends to infinity
the ratio of $R(y_n)/R(z_n)$ tends to infinity. This means that for
all $n$ sufficiently large, no point of  $\gamma_n$ can be contained
in an $\epsilon$-round component
or a $C$-component, because if it were then all
of $\gamma_n$ would be contained in that component, contradicting
the fact that the curvature ratio is arbitrarily large for large
$n$. Hence, for $n$ sufficiently large, every point of $\gamma_n$ is
either contained in the core of a
$(C,\epsilon)$-cap or is the center of a
strong $\epsilon$-neck. According to
Proposition~\ref{Xcontainedin}, for all $n$ sufficiently large
$\gamma_n$ is contained an open submanifold $X_n$ of the zero
time-slice of ${\mathcal M}'_n$ that is one of the following:
\begin{enumerate}
\item[(1)] an $\epsilon$-tube  and both endpoints of $\gamma_n$ are centers
of $\epsilon$-necks
contained in $X_n$,
\item[(2)] a $C$-capped $\epsilon$-tube with cap ${\mathcal C}$,
and each endpoint of $\gamma_n$ either is contained in the core $Y$
of ${\mathcal C}$ or is the center of an $\epsilon$-neck contained
in $X_n$,
\item[(3)] a double $C$-capped $\epsilon$-tube, or finally
\item[(4)] the union of two $(C,\epsilon)$-caps.
\end{enumerate}

The fourth possibility is incompatible with the fact that the ratio
of the curvatures at the endpoints of $\gamma_n$ grows arbitrarily
large as $n$ tends to infinity. Hence, this fourth possibility
cannot occur for $n$ sufficiently large. Thus, for all $n$
sufficiently large $X_n$ is one of the first three types listed
above.

\begin{claim}[Existence of a minimizing geodesic in $X_n$]
\label{claim:minimizing-geodesic-in-Xn}
There is a geodesic $\hat\gamma_n$ in $X_n$ with endpoints $z_n$ and
$y_n$. This geodesic is minimizing among all paths in $X_n$ from
$z_n$ to $y_n$.
\end{claim}

\begin{proof}
\sketch
This is clear in the third case since $X_n$ is a closed manifold.

Let us consider the first case. There are $\epsilon$-necks $N(z_n)$
and $N(y_n)$ centered at $z_n$ and $y_n$ and contained in $X_n$.
Suppose first that the central $2$-spheres $S(z_n)$ and $S(y_n)$ of
these necks are disjoint. Then they are the boundary of a compact
sumanifold $X_n'$ of $X_n$. It follows easily from
Lemma~\ref{curveshort} that any sequence of minimizing paths from
$z_n$ to $y_n$ is contained in the union of $X_n'$ with the  middle
halves of $N(z_n)$ and $N(y_n)$. Since this manifold has compact
closure in $X_n$, the usual arguments show that one can extract a
limit of a subsequence which is a minimizing geodesic in $X_n$ from
$z_n$ to $y_n$. If $S(z_n)\cap S(y_n)\not=\emptyset$, then $y_n$ is
contained in the middle half of $N(z_n)$, and again it follows
immediately from Lemma~\ref{curveshort} that there is a minimizing
geodesic in $N(z_n)$ between these points.

Now let us consider the second case. If each of $z_n$ and $y_n$ is
the center of an $\epsilon$-neck in $X_n$, the argument as in the
first case applies. If both points are contained in the core of
${\mathcal C}$ then, since that core has compact closure in $X_n$,
the result is again immediate. Lastly, suppose that one of the
points, we can assume by the symmetry of the roles of the points
that it is $z_n$, is the center of an $\epsilon$-neck $N(z_n)$ in
$X_n$ and the other is contained in the core of ${\mathcal C}$.
Suppose that the central $2$-sphere $S(z_n)$ of $N(z_n)$  meets the
core $Y$ of ${\mathcal C}$. Then $z_n$ lies in the  half of the neck
$N={\mathcal C}\setminus \overline Y$ whose closure contains the
frontier of $Y$. Orient $s_N$ so that this half is the positive
half. Thus, by Lemma~\ref{curveshort} any minimizing sequence of
paths from $z_n$ to $y_n$ is eventually contained in the union of
the core of ${\mathcal C}$ and the the positive three-quarters of
this neck. Hence, as before we can pass to a limit and construct  a
minimizing geodesic in $X_n$ connecting $z_n$ to $y_n$. On the other
hand, if $S(z_n)$ is disjoint from $Y$, then $S(z_n)$ separates
$X_n$ into a compact complementary component and a non-compact
complementary component and the compact complementary component
contains $Y$. Orient the $s_N$-direction so that the compact
complementary component lies on the positive side of $S(z_n)$. Then
any minimizing sequence of paths in $X_n$ from $z_n$ to $y_n$ is
eventually contained in the union of the compact complementary
component of $N(z_n)$ and the positive $3/4$'s of $N(z_n)$. As
before, this allows us to pass to a limit to obtain a minimizing
geodesic in $X_n$.
\end{proof}

This claim allows us to assume (as we now shall) that $\gamma_n$ is
a minimizing geodesic in $X_n$ from $z_n$ to $y_n$.


\begin{claim}[Balanced chain of necks along a sub-geodesic]
\label{claim:balanced-chain-subgeodesic}
For every $n$ sufficiently large, there is a sub-geodesic
$\gamma'_n$ of $\gamma_n$ with end points $z_n'$ and $y_n'$ such
that the following hold:
\begin{enumerate}
\item[(1)] The length of $\gamma'_n$ is bounded independent of $n$.
\item[(2)] $R(z_n')$ is bounded independent of $n$.
\item[(3)] $R(y_n')$ tends to infinity as $n$ tends to infinity.
\item[(4)] $\gamma_n'$ is contained in a strong $\epsilon$-tube
$T_n$ that is the union of a balanced
chain of strong
$\epsilon$-necks centered at points of $\gamma'_n$. The first
element in this chain is a strong $\epsilon$-neck $N(z'_n)$ centered
at $z'_n$. The last element is a strong $\epsilon$-neck containing
$y'_n$.
\item[(5)] For every  $x\in T_n$, we have $R(x)>3$ and $x$ is the center of a strong
$\epsilon$-neck in the flow $({\mathcal M}'_n,G'_n)$.
\end{enumerate}
\end{claim}

\begin{proof}
\sketch
The first item is clear since, for all $n$, the geodesic $\gamma_n$
has $G_n$-length at most $A_0Q_n^{-1/2}$ and hence $G'_n$-length at
most $A_0$. Suppose that we have a $(C,\epsilon)$-cap ${\mathcal C}$
whose core $Y$ contains a point of $\gamma_n$. Let $N$ be the
$\epsilon$-neck that is the complement of the closure of $Y$ in
${\mathcal C}$, and let $\widehat Y$ be the union of $Y$ and the
closed negative half of $N$. We claim that $\widehat Y$ contains
either $z_n$ or $y_n$. By Corollary~\ref{coreint}, since $Y$
contains a point of $\gamma_n$, the intersection of $\widehat Y$
with $\gamma_n$ is a subinterval containing one of the end points of
$\gamma_n$, i.e., either $z_n$ or $y_n$. This means that any point
$w$  which is contained in a $(C,\epsilon)$-cap whose core contains
a point of $\gamma_n$ must satisfy one of the following:
$$R(w)< CR(z'_n)\ \ \ \text{or}\ \ \ R(w)> C^{-1}R(y'_n).$$

We pass to a subsequence so that $R(y_n)/R(z_n)>4C^2$ for all $n$,
and we pass to a subinterval $\gamma'_n$ of $\gamma_n$ with
endpoints $z_n'$ and $y_n'$ such that:
\begin{enumerate}
\item[(1)] $R(z'_n)=2CR(z_n)$
\item[(2)] $R(y'_n)=(2C)^{-1}R(y_n)$
\item[(3)] $R(z'_n)\le R(w)\le R(y'_n)$ for all $w\in \gamma'_n$.
\end{enumerate}
Clearly, with these choices $R(z'_n)$ is bounded independent of $n$
and $R(y'_n)$ tends to infinity as $n$ tends to infinity. Also, no
point of $\gamma'_n$ is contained in the core of a
$(C,\epsilon)$-cap. Since every point of $\gamma'_n$ has a strong
$(C,\epsilon)$-canonical neighborhood, it follows that every point
of $\gamma'_n$ is the center of a strong $\epsilon$-neck. It now
follows from Proposition~\ref{epschains} that there is a balanced
$\epsilon$-chain consisting of strong $\epsilon$-necks centered at
points of $\gamma'_n$ whose union contains $\gamma'_n$. (Even if the
$2$-spheres of these necks do not separate the zero time-slice of
${\mathcal M}'_n$, as we build the balanced $\epsilon$-chain as
described in Proposition~\ref{epschains} the new necks we add can
not meet the negative end of $N(z'_n)$ since the geodesic
$\gamma'_n$ is minimal.) We can take the first element in the
balanced chain to be a strong $\epsilon$-neck $N(z'_n)$ centered at
$z'_n$, and the last element to be a strong $\epsilon$-neck $N^+_n$
containing $y'_n$. The union of this chain is $T_n$.

Next, we show that every point of $T_n$ is the center of a strong
$\epsilon$-neck in $({\mathcal M}_n,G_n)$. We must rule out the
possibility that there is a
 point of $T_n$ that is contained in the core of a $(C,\epsilon)$-cap.
Since $T_n$ is a union of $\epsilon$-necks centered at points of
$\gamma_n'$ we see that every point  $w\in T_n$ has
$$(3C/2)R(z_n)<R(w)<(2/3C)R(y_n).$$
This tells us that no point of $T_n$ is contained in a
$(C,\epsilon)$-cap whose core contains a point of $\gamma_n$. Thus,
to complete the argument we need only see that if there is a point
of $T_n$ contained in the core of a $(C,\epsilon)$-cap then the core
of that $(C,\epsilon)$-cap also contains a point of $\gamma_n$.
 The
scalar curvature inequality implies that both $z_n$ and $y_n$ are
outside $T_n$. This means that $\gamma_n$ traverses $T_n$ from one
end to the other. Let $w_-$, resp. $w_+$, be the point of $\gamma_n$
that lies in the frontier of $T_n$ contained in the closure of the
$N(z_n')$, resp. $N^+_n$. Since the scalar curvatures at these two
points of $\gamma$ satisfy the weak version of the above
inequalities, we see that there are strong $\epsilon$-necks $N(w_-)$
and $N(w_+)$ centered at them. Let $\widehat T_n$ be the union of
$T_n$, $N(w_-)$ and $N(w_+)$. It is also a strong $\epsilon$-tube,
and every point $\hat w$ of $\widehat T_n$ satisfies
$$(1.1)CR(z_n)<R(\hat w)<(0.9)C^{-1}R(y_n).$$
Thus, $z_n$ and $y_n$ are disjoint from $\widehat T_n$ and hence
$\gamma$ crosses $\widehat T_n$ from one end to the other.

 Now
suppose that $T_n$ meets the core $Y$ of a $(C,\epsilon)$-cap
${\mathcal C}$.
 Consider the boundary $S$ of the closure of $Y$.
If it is disjoint from $T_n$ then $T_n$ is contained in the core
$Y$. For large $n$ this is inconsistent with the fact that the ratio
of the scalar curvature at the endpoints of $\gamma'_n$ goes to
infinity. Thus, we are left to consider the case when
 $S$ contains a point of the tube $T_n$. In this case $S$ is
completely contained in $\widehat T_n$ and by
Corollary~\ref{s2isotopic} $S$ is isotopic to the $2$-spheres of the
product decomposition of $\widehat T_n$. Hence, $S$ meets a point of
$\gamma_n$ and consequently the core $Y$ contains a point of
$\gamma_n$. But we have already seen that this is not possible.

Lastly, we must show that $R(x)>3$ for every $x\in T_n$. We have
just seen that every $x\in T_n$ is the center of an $\epsilon$-neck.
If $x$ is contained in the $\epsilon$-neck centered at $z'_n$ or
$y'_n$, then since $R(z'_n)\ge 4$ and $R(y'_n)\ge 4$, clearly
$R(x)>3$. We must consider the case when $x$ is not contained in
either of these $\epsilon$-necks. In this case the central
$2$-sphere $S_x$of the $\epsilon$-neck centered at $x$ is contained
in the compact submanifold of $T_n'$ bounded by the central
$2$-spheres of the necks centered at $z'_n$ and $y'_n$. These
$2$-spheres are disjoint and by Condition 4 in
Proposition~\ref{S2intersection} each is a homotopically non-trivial
$2$-sphere in $T_n'$. Hence, the compact manifold with their
disjoint union as boundary is diffeomorphic to $S^2\times [0,1]$
and, again according to Condition 4 of
Proposition~\ref{S2intersection}, $S_x$ is isotopic to the
$2$-sphere factor in this product decomposition. Since the
intersection of $\gamma'_n$ with this submanifold is an arc spanning
from one boundary component to the other, $S_x$ must meet
$\gamma'_n$, in say $w$.
 By construction, since $w\in \gamma'_n$ we have $R(w)\ge 4$.
  This implies that $R(x)>3$.
 This completes the proof of the claim.
\end{proof}


\subsection{Extracting a limit of a subsequence of the tubes}

Passing to a subsequence we arrange that the $R(z_n')$ converge. Now
consider the subset ${\mathcal A}\subset \Ar$ consisting of all
$A>0$ such that there is a uniform bound, independent of $n$, for
the curvature on $B(z'_n,A)\cap T_n$. The set ${\mathcal A}$ is
non-empty since $R(z'_n)$ is bounded independent of $n$ and for
every $n$ there is a strong $\epsilon$-neck $N(z'_n)$ centered at
$z'_n$ contained in $T_n$. On the other hand, since
$d_{G_n'}(z_n',y_n')$ is uniformly bounded and $R(y_n')\rightarrow
\infty$,  there is a finite upper bound for ${\mathcal A}$. Let
$A_1$ be the least upper bound of ${\mathcal A}$. We set
$U_n=T_n\cap B(z'_n,A_1)$. This is an open subset of $T_n$
containing $z'_n$. We let $g_n'=G_n'|U_n$.

\begin{claim}[$U_n$ contains the neck centered at $z_n'$]
\label{claim:neck-contained-in-Un}
For all $n$ sufficiently large,  $3R(z'_n)^{-1/2}\epsilon^{-1}/2$ is
less than $A_1$, and hence $U_n$ contains the strong $\epsilon$-neck
$N(z_n')$ centered at $z'_n$.
\end{claim}

\begin{proof}
\sketch
The curvature on  $N(z'_n)$ is bounded independent of $n$. Consider
a point $w$ near the end of $N(z'_n)$ that separates $y_n'$ from
$z_n'$. It is also the center of a strong $\epsilon$-neck $N(w)$. By
Proposition~\ref{S2intersection} and our assumption that
$\epsilon\le \bar\epsilon(10^{-2})$, the scalar curvature on
$N(z'_n)\cup N(w)$ is between $(0.9)R(z'_n)$ and $(1.1)R(z'_n)$.
Since, by construction, the negative end of $N(z'_n)$ contains an
end of $T_n$, this implies that
$$N(z'_n)\cup N(w)\supset B(z'_n,7R(z'_n)^{-1/2}\epsilon^{-1}/4)\cap T_n,$$
 so that we see that
$A_1\ge 7\epsilon^{-1}\lim_{n\rightarrow\infty}
R(z'_n)^{-1/2}/4$. Thus, $A_1>3R(z'_n)^{-1}\epsilon^{-1}/2$ for all
$n$ sufficiently large. Obviously then $U_n$ contains $N(z'_n)$.
\end{proof}

The next claim uses terminology from Definition~\ref{deltareg}.

\begin{claim}[Uniform curvature bound on $\operatorname{Reg}_\delta(U_n,g_n')$]
\label{claim:uniform-curvature-bound-reg-delta}
For any $\delta>0$ there is a uniform bound, independent of $n$, for
the curvature on $\operatorname{Reg}_\delta(U_n,g'_n)$.
\end{claim}

\begin{proof}
\sketch
To prove this it suffices to show that given $\delta>0$ there is
$A<A_1$ such that $\operatorname{Reg}_\delta(U_n,g'_n)\subset B(z'_n,A)$ for
all $n$ sufficiently large. Of course, if we establish this for
every $\delta>0$ sufficiently small, then it follows for all
$\delta>0$. First of all, by Corollary~\ref{neckgeo} and
Lemma~\ref{neckcurv}, the fact that $\epsilon\le
\bar\epsilon(10^{-2})$ implies that any point $w$ with the property
that the strong $\epsilon$-neighborhood centered at $w$ contains
$z_n'$ is contained in the ball of radius
$(1.1)R(z'_n)^{-1/2}\epsilon^{-1}<A_1$ centered at $z_n'$. Thus, it
suffices to consider points $w_n$ in $\operatorname{Reg}_\delta(U_n,g'_n)$
with the property that the strong $\epsilon$-neck centered at $w_n$
does not contain $z_n'$. Fix such a $w_n$. Take a path $\mu_n(s)$
starting at $w_n$ moving in the $s$-direction at unit speed measured
in the $s$-coordinate of the $\epsilon$-neck centered at $w_n$ away
from $z'_n$ and ending at the frontier of this neck. Let $u_1$ be
the final point of this path. The rescaled version of
Lemma~\ref{neckdistcomp} implies that the forward difference
quotient for the distance from $z'_n$ satisfies
$$(0.99)R(w_n)^{-1/2}\le \frac{d}{ds}d(z'_n,\mu_n(s))
\le (1.01)R(w_n)^{-1/2}.$$ Of course, since we are working in an
$\epsilon$-neck we also have
$$(1-\epsilon)R(w_n)^{-1/2}\le \frac{d(d(w_n,\mu_n(s)))}{ds}\le
(1+\epsilon)R(w_n)^{-1/2}.$$  We continue the path $\mu_n$ moving in
the $s$-direction of a neck centered at $u_1$. Applying
Lemma~\ref{neckdistcomp} again both to the distance from $w_n$ and
the distance from $z'_n$ yields:
$$(0.99)R(u_1)^{-1/2}\le \frac{d(d(z'_n,\mu_n(s)))}
{ds}\le (1.01)R(u_1)^{-1/2}$$
$$(0.99)R(u_1)^{-1/2}\le \frac{d(d(w_n,\mu_n(s)))}
{ds}\le (1.01)R(u_1)^{-1/2}$$ on this part of the path $\mu_n$. We
repeat this process as many times as necessary until we reach a
point $w'_n\in U_n$ at distance $\delta/2$ from $w_n$. This is
possible since the ball of radius $\delta$ centered at $w_n$ is
contained in $U_n$. By the difference quotient inequalities, it
follows that $d(z'_n,w'_n)-d(z'_n,w_n)>\delta/4$.
 Since $w_n'\in U_n$ and
consequently that $d(z_n',w_n')<A_1$. It follows that
$d(z_n',w_n)\le A_1-\delta/4$. This proves that, for all $n$
sufficiently large, $\operatorname{Reg}_\delta(U_n,g'_n)\subset
B(z_n',A_1-\delta/4)$, and consequently that the curvature on
$\operatorname{Reg}_\delta(U_n,g'_n)$ is bounded independent of $n$.
\end{proof}

By Shi's theorem (Theorem~\ref{thm:shi-derivative-estimates}), the fact that each point of
$U_n$ is the center of a strong $\epsilon$-neck means that there is
a bound, independent of $n$, on all covariant derivatives of the
curvature at any point of $U_n$ in terms of the bound on the
curvature at the center point. In particular, because of the
previous result, we see that for any $\epsilon>0$ and any $\ell\ge
0$ there is a uniform bound for $|\nabla^\ell \Rm|$ on
$\operatorname{Reg}_\delta(U_n,g'_n)$. Clearly, since the base point $z_n'$ has
bounded curvature it lies in $\operatorname{Reg}_\delta(U_n,g'_n)$ for
sufficiently small $\delta$ (how small being independent of $n$).
Lastly, the fact that every point in $U_n$ is the center of an
$\epsilon$-neighborhood implies that $(U_n,g'_n)$ is $\kappa$
non-collapsed on scales $\le r_0$ where both $\kappa$ and $r_0$ are
universal. Since the $\gamma'_n$ have uniformly bounded lengths, the
$\epsilon$-tubes $T_n'$ have uniformly bounded diameter. Also, we
have seen that their have curvatures are bounded from below by $3$.
It follows that their volumes are uniformly bounded. Now invoking
Theorem~\ref{basicconv} we see that after passing to a subsequence
we have a geometric limit $(U_\infty,g_\infty,z_\infty)$ of a
subsequence of $(U_n,g'_n,z'_n)$.


\subsection{Properties of the limiting tube}


Now we come to a result establishing all the properties we need for
the limiting manifold.


\begin{proposition}[Properties of the incomplete geometric limit tube]
\label{prop:incomplete-limit-tube-properties}
\uses{def:generalized-ricci-flow, def:evolving-epsilon-neck, def:epsilon-neck, claim:balanced-chain-subgeodesic, claim:uniform-curvature-bound-reg-delta, basicconv, thm:shi-derivative-estimates, smoothconv, canonvary, Xcontainedin, partialflowlimit, blowupposcurv}
\label{Uinfty}

The geometric limit $(U_\infty,g_\infty,z_\infty)$ is an incomplete
Riemannian $3$-manifold of finite diameter. There is a
diffeomorphism $\psi\colon U_\infty\to S^2\times (0,1)$. There is a
$2\epsilon$-neck centered at $z_\infty$ whose central $2$-sphere
$S^2(z_\infty)$ maps under $\psi$ to a $2$-sphere isotopic to a
$2$-sphere  factor in the product decomposition. The scalar
curvature is bounded at one end of $U_\infty$ but tends to infinity
at the other end, the latter end  which is denoted ${\mathcal E}$.
Let ${\mathcal U}_\infty\subset U_\infty\times(-\infty,0]$ be the
open subset consisting of all $(x,t)$ for which $-R(x)^{-1}<t\le 0$.
We have a generalized Ricci flow on ${\mathcal U}_\infty$ which is a
partial geometric limit of a subsequence of the generalized Ricci
flows $({\mathcal M}'_n,G'_n,z'_n)$. In particular, the zero-time
slice of the limit flow is $(U_\infty,g_\infty)$. The Riemannian
curvature is non-negative at all points of the limiting smooth flow
on ${\mathcal U}_\infty$. Every point $x\in U_\infty\times\{0\}$
which is not separated from ${\mathcal E}$ by $S^2(z_\infty)$ is the
center of an evolving $2\epsilon$-neck $N(x)$ defined for an
interval of normalized time of length $1/2$. Furthermore, the
central $2$-sphere of $N(x)$ is isotopic to the $2$-sphere factor of
$U_\infty$ under the diffeomorphism $\psi$.
\end{proposition}


The proof of this proposition occupies the rest of
Chapter~\ref{sectincomp}.

\begin{claim}[Points of $W^+_\infty$ are neck centers]
\label{claim:points-in-Wplus-are-neck-centers}
Any point $q\in W^+_\infty$ is the center of a $2\epsilon$-neck in
$U_\infty$.
\end{claim}
\begin{proof}
\sketch The approximating necks have uniform regular curvature bounds,
so smooth convergence produces a centered $2\epsilon$-neck at $q$.
\end{proof}

\begin{claim}[Curvature bounded at one end, unbounded at the other]
\label{claim:curvature-bounded-one-end-infinite-other}
The curvature is bounded in a neighborhood of one end of $U_\infty$
and goes to infinity at the other end.
\end{claim}
\begin{proof}
\sketch The bounded end is the limit of negative neck halves, while
neck scales tend to zero along the other end.
\end{proof}

\begin{claim}[Existence of a backward flowline]
\label{claim:flowline-exists-for-x}
For each $x\in U_n\subset M_n$ there is a flowline $\{x\}\times
(-R(x)^{-1},0]$ in ${\mathcal M}_n$, and its scalar curvature is at
most $R(x)$.
\end{claim}
\begin{proof}
\sketch A strong $\epsilon$-neck centered at $x$ supplies the flowline
and the curvature bound.
\end{proof}

\begin{claim}[Non-negativity of the limiting curvature]
\label{claim:limiting-flow-nonnegative-curvature}
The curvature of the generalized Ricci flow on ${\mathcal
U}_\infty$ is non-negative.
\end{claim}
\begin{proof}
\sketch Rescaling Equation~\ref{weakerpos} and using $Q_n\to\infty$
gives non-negative limiting curvature by Theorem~\ref{blowupposcurv}.
\end{proof}

\begin{proof}
\sketch
Let $V_1\subset V_2 \subset \cdots \subset U_\infty$ be the open
subsets and $\varphi_n\colon V_n\to U_n$ be the maps having all the
properties stated in Definition~\ref{smoothconv} so as to exhibit
$(U_\infty,g_\infty,z_\infty)$ as the  geometric limit of the
$(U_n,g_n',z_n')$.

Since the $U_n$ are all contained in $B(z_n',A_1)$, it follows that any point
of $U_\infty$ is within $A_1$ of the limiting base point $z_\infty$. This
proves that the diameter of $U_\infty$ is bounded.




For each $n$ there is the  $\epsilon$-neck $N(z_n')$ centered at
$z_n'$ contained in $U_n$. The middle two-thirds, $N_n'$, of this
neck has closure contained in $\operatorname{Reg}_\delta(U_n,g_n)$ for some
$\delta>0$ independent of $n$ (in fact, restricting to $n$
sufficiently large, $\delta$ can be taken to be approximately equal
to $R(z_\infty)^{-1/2}\epsilon^{-1}/3$). This means that for some
$n$ sufficiently large and for all $m\ge n$ the image
$\varphi_m(V_n)\subset U_m$ contains $N'_m$. For any fixed $n$ as
$m$ tends to infinity  the metrics $\varphi_m^*g_m|_{V_n}$ converge
uniformly in the $C^\infty$-topology to $g_\infty|_{V_n}$. Thus, it
follows from Proposition~\ref{canonvary} that for all $m$
sufficiently large, $\varphi^{-1}_m(N'_m)$ is a $3\epsilon/2$-neck
centered at $z_\infty$. We fix such a neck $N'(z_\infty)\subset
U_\infty$. Let $S(z_\infty)$ be the central $2$-sphere of
$N'(z_\infty)$. For each $n$ sufficiently large,
$\varphi_n(S(z_\infty))$ separates $U_n$ into two components, one,
say $W^-_n$ contained in $N(z'_n)$ and the other, $W^+_n$ containing
all of $U_n\setminus N(z'_n)$. It follows that $S(z_\infty)$
separates $U_\infty$ into two components, one, denoted $W^-_\infty$,
where the curvature is bounded (and where, in fact, the curvature is
close to $R(z_\infty)$) and the other, denoted $W^+_\infty$, where
it is unbounded.

Fix a point $q\in W^+_\infty$.
 For all $n$ sufficiently large denote by $q_n=\varphi_n(q)$.
 Then for all $n$ sufficiently large, $q_n\in W^+_n$ and
$\lim_{n\rightarrow\infty}R(q_n)=R(q)$. This means that for all
$n$ sufficiently large $R(y_n')>>R(q_n))$, and hence
 the $3\epsilon/2$-neck centered at
 $q_n\in U_n$ is disjoint from $N(y'_n)$. Thus, by the rescaled version of
 Corollary~\ref{neckgeo}, we see that  the distance from
 the $3\epsilon/2$-neck centered at $q_n$ to $N(y'_n)$ is bounded below by
 $(0.99)\epsilon^{-1}R(q_n)^{-1/2}/4\ge \epsilon^{-1}R(q_\infty)^{-1/2}/12$.
 Also, since $q_n\in W_n$, this $3\epsilon/2$-neck $N'(q_n)$ centered at
 $q_n$ does not extend past the $2$-sphere at $s^{-1}(-3\epsilon^{-1}/4)$ in the
 $ \epsilon$-neck $N(z'_n)$. It follows that for all $n$ sufficiently large that
 this $3\epsilon/2$-neck has
 compact closure contained in $\operatorname{Reg}_\delta(U_n,g_n)$ for some $\delta$
 independent of $n$, and hence there is $m$ such that for all $n$ sufficiently large $N'(q_n)$ is
 contained in the image $\varphi_n(V_m)$. Again using the fact that
 $\varphi_n^*(g_n|_{V_m})$ converges in the $C^\infty$-topology to
 $g_\infty|_{V_m}$ as $n$ tends to infinity, we see, by Proposition~\ref{canonvary}
 that for all $n$
 sufficiently large $\varphi_n^{-1}(N_m)$ contains a $2\epsilon$-neck in $U_\infty$ centered
 at $q$.


It now follows from Proposition~\ref{Xcontainedin} that $W^+_\infty$
is contained in an $2\epsilon$-tube $T_\infty$ that is contained in
$U_\infty$. Furthermore, the frontier of $W^+_\infty$ in $T_\infty$
is the $2$-sphere $S(z_\infty)$ which is isotopic to the central
$2$-spheres of the $2\epsilon$-necks making up $T_\infty$. Hence,
the closure $\overline W^+_\infty$ of $W^+_\infty$ is a
$2\epsilon$-tube with boundary $S(z_\infty)$. In particular,
$\overline W^+_\infty$ is diffeomorphic to $S^2\times [0,1)$.

Now we consider the closure $\overline W^-_\infty$ of  $W^-_\infty$.
Since the closure of each $W^-_n$ is the closed negative half of the
$\epsilon$-neck $N(z_n')$ and the curvatures of the $z_n'$ have a
finite, positive limit, the limit $\overline W^-_\infty$ is
diffeomorphic to a product $S^2\times (-1,0]$. Hence, $U_\infty$ is
the union of  $\overline W^+_\infty$ and  $\overline W^-_\infty$
along their common boundary. It follows immediately that $U_\infty$
is diffeomorphic to $S^2\times (0,1)$.


 A neighborhood of one end of $U_\infty$, the end $\overline W^-_\infty$,
 is the limit of the
 negative halves of $\epsilon$-necks centered at $z_n'$.
 Thus, the curvature is bounded on this neighborhood, and in fact is
approximately equal to $R(z_\infty)$. Let $x_k$ be any sequence of
points in $U_\infty$ tending to the other end. We show that $R(x_k)$
tends to $\infty$ as $k$ does. The point is that since the sequence
is tending to the end, the distance from $x_k$ to the end of
$U_\infty$ is going to zero. Yet, each $x_k$ is the center of an
$\epsilon$-neck in $U_\infty$. The only way this is possible is if
the scales of these $\epsilon$-necks are converging to zero as $k$
goes to infinity. This is equivalent to the statement that $R(x_k)$
tends to $\infty$ as $k$ goes to infinity.


The next step in the proof of Proposition~\ref{prop:incomplete-limit-tube-properties} is to extend
the flow backwards a certain amount. As stated in the proposition,
the amount of backward time that we can extend the flow is not
uniform over all of $U_\infty$, but rather depends on the curvature
of the point at time zero.



Since $x\in U_n\subset T_n$, there is a strong $\epsilon$-neck in
${\mathcal M}_n$ centered at $x$. Both statements follow immediately
from that.


Let $X\subset U_\infty$ be an open submanifold with compact closure
and set
$$t_0(X)= \sup_{x\in X}(-R_{g_\infty}(x)^{-1}).$$ Then for
all $n$ sufficiently large $\varphi_n$ is defined on $X$ and the
scalar curvature of the flow $g_n(t)$ on $\varphi_n(X)\times
(t_0,0]$ is uniformly bounded independent of $n$. Thus, according to
Proposition~\ref{partialflowlimit} by passing to a subsequence we
can arrange that there is a limiting flow defined on $X\times
(t_0,0]$. Let ${\mathcal U}_\infty\subset U_\infty\times
(-\infty,0]$ consist of all pairs $(x,t)$ with the property that
$-R_{g_\infty}(x,0)^{-1}<t\le 0$.
 Cover ${\mathcal U}_\infty$ by countably many such boxes
of the type $X\times (-t_0(X),0]$ as described above, and take a
diagonal subsequence. This allows us to pass to a subsequence so
that the limiting flow exists (as a generalized Ricci flow) on
${\mathcal U}_\infty$.



 This claim follows from the fact that the
original sequence $({\mathcal M}_n,G_n)$ consists of generalized
flows whose curvatures are pinched toward positive in the weak sense
given in Equation~\ref{weakerpos} and the fact that $Q_n\rightarrow
\infty$ as $n\rightarrow \infty$. (See Theorem~\ref{blowupposcurv}.)

\end{proof}





\section{Cone limits near the end ${\mathcal E}$ for rescalings of $U_\infty$}\label{GromHaus}

We study rescalings of $U_\infty$ near ${\mathcal E}$.

Let $(X,d_X)$ be a metric space. Recall that the  cone on $X$,
denoted $C(X)$, is the quotient space $X\times [0,\infty)$ under the
identification $(x,0)\cong (y,0)$ for all $x,y\in X$. The image of
$X\times\{0\}$ is the {\sl cone point} of the cone. The metric on
$C(X)$ is given by
\begin{equation}\label{conemetric}
d((x,s_1),(y,s_2))=s_1^2+s_2^2-2s_1s_2\cos(\mathrm{min}(d_X(x,y),\pi)). \end{equation} The open cone $C'(X)$ is the
complement of the cone point in $C(X)$ with the induced metric.


The purpose of this section is to prove the following result.

\begin{proposition}[Rescaled limits converge to a cone at the end]
\label{prop:cone-limit-gromov-hausdorff}
\uses{prop:incomplete-limit-tube-properties}
\label{statement}

Let $(U_\infty,g_\infty,z_\infty)$  be as in the conclusion of
Proposition~\ref{prop:incomplete-limit-tube-properties}. Let $Q_\infty=R_{g_\infty}(z_\infty)$ and
let ${\mathcal E}$ be the end of $U_\infty$ where the scalar
curvature is unbounded. Let $\lambda_n$ be any sequence of positive
numbers with $\lim_{n\rightarrow\infty}\lambda_n=+\infty$. Then
there is a sequence $x_n$ in $U_\infty$ such that for each $n$ the
distance from $x_n$ to ${\mathcal E}$ is $\lambda_n^{-1/2}$, and
such that the pointed Riemannian manifolds
$(U_\infty,\lambda_ng_\infty,x_n)$ converge in the Gromov-Hausdorff
sense to an open cone, an open cone not homeomorphic to
an open ray (i.e., not homeomorphic to the open cone on a point).
\end{proposition}

The rest of this section is devoted to the proof of this result.

\subsection{Directions at ${\mathcal E}$}

 We orient the direction down the tube $U_\infty$ so that
${\mathcal E}$ is at the positive end. This gives an $s_N$-direction
for each $2\epsilon$-neck $N$ contained in $U_\infty$.




Fix a point $x\in U_\infty$. We say a ray $\gamma$ with endpoint $x$
limiting to ${\mathcal E}$ is {\sl a minimizing geodesic ray} if for
every $y\in \gamma$ the segment on $\gamma$ from $x$ to $y$ is a
minimizing geodesic segment; i.e., the length of this geodesic
segment is equal to $d(x,y)$.

\begin{claim}[Existence of minimizing rays to the end]
\label{claim:minimizing-ray-to-end-exists}
\uses{prop:incomplete-limit-tube-properties, genint}

There is a minimizing geodesic ray to ${\mathcal E}$ from each $x\in
U_\infty$ with $R(x)\ge 2Q_\infty$.
\end{claim}

\begin{proof}
\sketch
Fix $x$ with $R(x)\ge 2Q_\infty$ and fix a $2\epsilon$-neck $N_x$
centered at $x$. Let $S_x^2$ be the central $2$-sphere of this neck.
Take a sequence of points $q_n$ tending to the end ${\mathcal E}$,
each being closer to the end than $x$ in the sense that $S^2_x$ does
not separate any $q_n$ from the end ${\mathcal E}$. We claim that
there is a minimizing geodesic from $x$ to each $q_n$. The reason is
that by Lemma~\ref{genint} any minimizing sequence of arcs from $x$
to $q_n$ cannot exit from the minus end of $N_x$ nor the plus end of
a $2\epsilon$-neck centered at $q_n$. Consider a sequence of paths
from $x$ to $q_n$ minimizing the distance.  Hence these paths all
lie in a fixed compact subset of $U_n$. After replacing the sequence
by a subsequence, we can pass to a limit, which is clearly a
minimizing geodesic from $x$ to $q_n$. Consider minimizing geodesics
$\mu_n$ from $x$ to $q_n$. The same argument shows that, after
passing to a subsequence, the $\mu_n$ converge to a minimizing
geodesic ray from $x$ to ${\mathcal E}$.
\end{proof}


\begin{claim}[Minimizing rays realize the distance to the end]
\label{claim:shortest-ray-equivalence}
(1) Any minimizing geodesic ray from $x$ to the end ${\mathcal E}$
is a shortest ray from $x$ to the end ${\mathcal E}$, and conversely
any shortest ray from $x$ to the end ${\mathcal E}$ is a minimizing
geodesic ray.

(2) The length of a shortest ray from $x$ to ${\mathcal E}$ is the
distance (see Section~\ref{ENDS}) from $x$ to ${\mathcal E}$.
\end{claim}


\begin{proof}
\sketch
The implication in (1) in one direction is clear: If $\gamma$ is a
ray from $x$ to the end ${\mathcal E}$, and for some $y\in \gamma$
the segment on $\gamma$ from $x$ to $y$ is not minimizing, then
there is a shorter geodesic segment $\mu$ from $x$ to $y$. The union
of this together with the ray on $\gamma$ from  $y$ to the end is a
shorter ray from $x$ to the end.

Let us establish the opposite implication. Suppose that $\gamma$ is
a minimizing geodesic ray from $x$ to the end ${\mathcal E}$ and
that there is a $\delta>0$ and a shortest geodesic ray $\gamma'$
from $x$ to the end ${\mathcal E}$ with $|\gamma'|=|\gamma|-\delta$.
As we have just seen, $\gamma'$ is a minimizing geodesic ray. Take a
sequence of points $q_i$ tending to the end ${\mathcal E}$ and let
$S^2_i$ be the central $2$-sphere in the $2\epsilon$-neck centered
at $q_i$. Of course, for all $i$ sufficiently large, both $\gamma'$
and $\gamma$ must cross $S^2_i$. Since the scalar curvature tends to
infinity at the end ${\mathcal E}$, it follows from
Lemma~\ref{directions} for all $i$ sufficiently large, the extrinsic
diameter of $S^2_i$ is less than $\delta/3$. Let $p_i$ be a point of
intersection of $\gamma$ with $S^2_i$. For all $i$ sufficiently
large the length $d_i$ of the sub-ray in $\gamma$ from $p_i$ to the
end ${\mathcal E}$ is at most $\delta/3$. Let $p'_i$ be a point of
intersection of $\gamma'$ with $S^2_i$ and let $d'_i$ be the length
of the ray in $\gamma'$ from $p'_i$ to the end ${\mathcal E}$. Let
$\lambda$ be the sub-geodesic of $\gamma$ from $x$ to $p_i$ and
$\lambda'$ the sub-geodesic of $\gamma'$ from $x$ to $p'_i$. Let
$\beta$ be a minimizing geodesic from $p'_i$ to $p_i$. Of course,
$|\beta|<\delta/3$ so that by the minimality of $\lambda$ and
$\lambda'$ we have
$$-\delta/3<|\lambda|-|\lambda'|<\delta/3.$$ Since
$|\lambda'|+d'_i=|\lambda|+d_i-\delta$, we have
$$2\delta/3\le d_i-d'_i.$$
This is absurd since $d'_i>0$ and $d_i<\delta/3$.

(2) follows immediately from (1) and the definition.
\end{proof}

Given this result, the usual arguments show:

\begin{corollary}[Uniqueness of shortest continuations]
\label{cor:unique-shortest-subray}
\uses{claim:shortest-ray-equivalence}
\label{distlim}

If $\gamma$ is a minimizing geodesic ray from $x$ to the end
${\mathcal E}$, then for any $y\in \gamma\setminus\{x\}$ the sub-ray
of $\gamma$ from $y$ to the end, is the unique shortest geodesic
from $y$ to the end.
\end{corollary}

Also, we have a version of the triangle inequality for distances to
${\mathcal E}$.

\begin{lemma}[Triangle inequality for distance to the end]
\label{lem:triangle-inequality-distance-to-end}
Let $x$ and $y$ be points of $M$. Then the three distances $d(x,y)$,
$d(x,{\mathcal E})$ and $d(y,{\mathcal E})$ satisfy the triangle
inequality.
\end{lemma}


\begin{proof}
\sketch
From the definitions it is clear that $d(x,y)+d(y,{\mathcal E})\ge
d(x,{\mathcal E})$, and symmetrically, reversing the roles of $x$
and $y$. The remaining inequality that we must establish is the
following: $d(x,{\mathcal E})+d(y,{\mathcal E})\ge d(x,y)$. Let
$q_n$ be any sequence of points converging to ${\mathcal E}$. Since
the end is at finite distance, it is clear that $d(x,{\mathcal E})=
\lim_{n\rightarrow\infty}d(x,q_n)$. The remaining inequality
follows from this and  the usual triangle inequality applied to
$d(x,q_n)$, $d(y,q_n)$ and $d(x,y)$.
\end{proof}


\begin{definition}[Direction at the end]
\label{def:direction-at-end}
\uses{cor:unique-shortest-subray}

We say that two minimizing geodesic rays limiting to ${\mathcal E}$
are {\em equivalent} if one is contained in the other. From the
unique continuation of geodesics it is easy to see that this
generates an equivalence relation. An equivalence class is a {\sl
direction at ${\mathcal E}$}, and the set of equivalence classes is
the {\sl set of directions at ${\mathcal E}$}.
\end{definition}

\begin{lemma}[Multiple directions at the end]
\label{lem:multiple-directions-at-end}
\uses{def:direction-at-end}
\label{morethanone}

There is more than one direction at ${\mathcal E}$.
\end{lemma}

\begin{proof}
\sketch
Take a minimal geodesic ray $\gamma$ from a point $x$ limiting to
the end and let $y$ be a point closer to ${\mathcal E}$ than $x$ and
not lying on $\gamma$. Then a minimal geodesic ray from $y$ to
${\mathcal E}$ gives a direction at ${\mathcal E}$ distinct from the
direction determined by $\gamma$.
\end{proof}

\begin{remark}[Space of directions is homeomorphic to $S^2$]
\label{rem:space-of-directions-topology}
In fact, the general theory of positively curved spaces implies that
the space of directions is homeomorphic to $S^2$. Since we do not
need this stronger result we do not prove it.
\end{remark}

\subsection{The Metric on the space of directions at ${\mathcal
E}$}



\begin{definition}[Comparison angle between two directions]
\label{def:angle-between-directions-theta}
\uses{def:direction-at-end}
\label{defntheta}

Let $\gamma$ and $\mu$ be minimizing geodesic rays limiting to
${\mathcal E}$, of lengths $a$ and $b$, parameterized by the
distance from the end. For $0<s\le a$ and $0<s'\le b$ construct a
triangle $\alpha_se\beta_{s'}$ in the Euclidean plane with
$|\alpha_se|=s, |e\beta_{s'}|=s'$ and
$|\alpha_s\beta_{s'}|=d(\gamma(s),\mu(s'))$. We define
$\theta(\gamma,s,\mu,s')$ to be the angle at $e$ of the triangle
$\alpha_se\beta_{s'}$.
\end{definition}

\begin{lemma}[Monotone convergence of comparison angles]
\label{lem:theta-monotone-convergence}
\uses{def:angle-between-directions-theta, thm:length-comparison}
\label{thetaconv}

For all $\gamma,s,\mu,s'$ as in the previous definition we have
$$0\le \theta(\gamma,s,\mu,s')\le \pi.$$ Furthermore,  $\theta(\gamma,s,\mu,s')$
is a non-increasing function of $s$  when $\gamma,\mu,s'$ are held
fixed, and symmetrically it is a non-increasing function of $s'$
when $\gamma,s,\mu$ are held fixed. In particular, fixing $\gamma$
and $\mu$, the function $\theta(\gamma,s,\mu,s')$ is non-decreasing
as $s$ and $s'$ tend to zero. Thus, there is a well-defined limit as
$s$ and $s'$ go to zero, denoted $ \theta(\gamma,\mu)$. This limit
is greater than or equal to $\theta(\gamma,s,\mu,s')$ for all $s$
and $s'$ for which the latter is defined. We have
$0\le\theta(\gamma,\mu)\le \pi$. The angle $\theta(\gamma,\mu)=0$ if
and only if $\gamma$ and $\mu$ are equivalent. Furthermore, if
$\gamma$ is equivalent to $\gamma'$ and $\mu$ is equivalent to
$\mu'$, then $\theta(\gamma,\mu)=\theta(\gamma',\mu')$.
\end{lemma}


\begin{proof}
\sketch
By restricting $\gamma$ and $\mu$ to slightly smaller rays, we can
assume that each is the unique shortest ray from its endpoint to the
end ${\mathcal E}$. Let $x$, resp., $y$ be the endpoint of $\gamma$,
resp., $\mu$. Now let $q_n$ be any sequence of points in $U_\infty$
limiting to the end ${\mathcal E}$, and consider minimizing geodesic
rays $\gamma_n$ from $q_n$ to $x$ and $\mu_n$ from $q_n$ to $y$,
each parameterized by the distance from $q_n$. By passing to a
subsequence we can assume that each of the sequences $\{\gamma_n\}$
and $\{\mu_n\}$ converge to a minimizing geodesic ray, which  by
uniqueness, implies that the first sequence limits to $\gamma$ and
the second to $\mu$. For $s,s'$ sufficiently small, let
$\theta_n(s,s')$ be the angle at $\tilde q_n$ of the Euclidean
triangle $\alpha_n\tilde q_n\beta_n$, where $|\alpha_n\tilde
q_n|=d(\gamma_n(s),q_n)$, $|\beta_n\tilde q_n|=d(\mu_n(s'),q_n)$ and
$|\alpha_n\beta_n|=d(\gamma_n(s),\mu_n(s'))$. Clearly, for fixed $s$
and $s'$ sufficiently small, $\theta_n(s,s')$ converges as
$n\rightarrow\infty$ to $\theta(\gamma,s,\mu,s')$. By the Toponogov
property (Theorem~\ref{thm:length-comparison}) for manifolds with
non-negative curvature, for each $n$ the function $\theta_n(s,s')$
is a non-increasing function of each variable, when the other is
held fixed. This property then passes to the limit, giving the first
statement in the lemma.

By the monotonicity, $\theta(\gamma,\mu)=0$ if and only if for all
$s,s'$ sufficiently small we have $\theta(\gamma,s,\mu,s')=0$, which
means one of $\gamma$ and $\mu$ is contained in the other.

It is obvious that the last statement holds.
\end{proof}


It follows that $\theta(\gamma,\mu)$ yields a well-defined function
on the set of pairs of directions at ${\mathcal E}$. It is clearly a
symmetric, non-negative function which is positive off of the
diagonal. The next lemma shows that it is a metric by establishing
the triangle inequality for $\theta$.

\begin{lemma}[Triangle inequality for the angle metric]
\label{lem:theta-triangle-inequality}
\uses{def:angle-between-directions-theta}

If $\gamma,\mu,\nu$ are minimizing geodesic rays limiting to
${\mathcal E}$, then
$$\theta(\gamma,\mu)+\theta(\mu,\nu)\ge \theta(\gamma,\nu).$$
\end{lemma}


\begin{proof}
\sketch
\uses{cor:unique-shortest-subray,cor:angle-comparison}
By Corollary~\ref{cor:unique-shortest-subray}, after replacing
$\gamma, \mu,\nu$ by
equivalent, shorter geodesic arcs, we can assume that they are the
unique minimizing geodesics from their end points, say $x,y,z$
respectively, to ${\mathcal E}$. Let $q_n$ be a sequence of points
limiting to ${\mathcal E}$, and let $\gamma_n,\mu_n,\nu_n$ be
minimizing geodesics from $x,y,z$ to $q_n$. Denote by
$\theta_n(x,y), \theta_n(y,z)$, and $\theta_n(x,z)$, respectively,
the angles  at $\tilde q_n$ of the Euclidean triangles with the
following edge lengths: $\{d(x,y), d(x,q_n), d(y,q_n)\}$,
$\{d(y,z),d(y,q_n),d(z,q_n)\}$, and $\{d(z,x),d(z,q_n),d(x,q_n)\}$.
According to Corollary~\ref{cor:angle-comparison} we have
$\theta_n(x,y)+\theta_n(y,z)\ge \theta_n(x,z)$. Passing to the limit
as $n$ goes to $\infty$ and then the limit as $x$, $y$ and $z$ tend
to ${\mathcal E}$, gives the result.
\end{proof}


\begin{definition}[Cone of directions at the end]
\label{def:cone-of-directions}
\uses{def:angle-between-directions-theta, def:direction-at-end}

Let $X({\mathcal E})$ denote the set of directions at ${\mathcal
E}$. We define the metric on $X({\mathcal E})$ by setting
$d([\gamma],[\mu])=\theta(\gamma,\mu)$. We call this the {\sl
(metric) space of realized directions at} ${\mathcal E}$. The {\sl
metric space of directions at} ${\mathcal E}$ is the completion
$\bar X({\mathcal E})$ of $X({\mathcal E})$ with respect to the
given metric. We denote by $(C_{\mathcal E},g_{\mathcal E})$ the
cone on $\bar X({\mathcal E})$ with the cone metric as given in
Equation~(\ref{conemetric}).
\end{definition}


\begin{proposition}[The cone of directions is not a ray]
\label{prop:cone-not-a-ray}
\uses{def:cone-of-directions, lem:multiple-directions-at-end}
\label{notray}

$(C_{\mathcal E},g_{\mathcal E})$ is a metric cone that is not
homeomorphic to a ray.
\end{proposition}


\begin{proof}
\sketch
\uses{lem:multiple-directions-at-end}
By construction $(C_{\mathcal E},g_{\mathcal E})$ is a metric cone.
That it is not homeomorphic to a ray follows immediately from
Lemma~\ref{lem:multiple-directions-at-end}.
\end{proof}


\subsection{Comparison results for distances}

\begin{lemma}[Cone distance bounds the limit distance]
\label{lem:cone-distance-upper-bound}
\uses{def:cone-of-directions}
\label{distcomp}

Suppose that $\gamma$ and $\mu$ are unique shortest geodesic rays
from points $x$ and $y$ to the end ${\mathcal E}$. Let $[\gamma]$
and $[\mu]$ be the points of $X({\mathcal E})$ represented by these
two geodesics rays. Let $a$, resp. $b$, be the distance from $x$,
resp. $y$, to ${\mathcal E}$. Denote by  $x'$, resp. $y'$, the image
in $C_{\mathcal E}$ of the point $([\gamma],a)$, resp. $([\mu],b)$,
of $X({\mathcal E})\times [0,\infty)$. Then
$$d_{g_\infty}(x,y)\le d_{g_{\mathcal E}}(x',y').$$
\end{lemma}

\begin{proof}
\sketch
\uses{def:angle-between-directions-theta,lem:theta-monotone-convergence}
By the definition of the cone metric we have
$$d_{g_{\mathcal E}}(x',y')=a^2+b^2-2ab\,\cos(\theta(\gamma,\mu)).$$
On the other hand by Definition~\ref{def:angle-between-directions-theta} and the law of
cosines for Euclidean triangles, we have
$$d_{g_\infty}(x,y)=a^2+b^2-2ab\,\cos(\theta(\gamma,a,\mu,b)).$$
The result is now immediate from the fact, proved in
Lemma~\ref{lem:theta-monotone-convergence} that
$$0\le \theta(\gamma,a,\mu,b)\le \theta(\gamma,\mu)\le \pi,$$
and the fact that the cosine is a monotone decreasing function on
the interval $[0,\pi]$.
\end{proof}

\begin{corollary}[Rescaled cone distance bound]
\label{cor:rescaled-cone-distance-bound}
\uses{lem:cone-distance-upper-bound}
\label{refdistcomp}

Let $\gamma,\mu,x,y$ be as in the previous lemma. Fix $\lambda>0$.
Let $a=d_{\lambda g_\infty}(x,{\mathcal E})$ and $b=d_{\lambda
g_\infty}(y,{\mathcal E})$. Set $x'_\lambda$ and $y'_\lambda$ equal
to the points in the cone $([\gamma],a)$ and $([\mu],b)$. Then we
have
$$d_{\lambda g_\infty}(x,y)\le d_{g_{\mathcal E}}(x'_\lambda,y'_\lambda).$$
\end{corollary}


\begin{proof}
\sketch
\uses{lem:cone-distance-upper-bound}
This is immediate by applying the previous lemma to the rescaled
manifold $(U_\infty,\lambda g_\infty)$, and noticing that rescaling
does not affect the cone $C_{\mathcal E}$ nor its metric.
\end{proof}

\begin{lemma}[Nets of directions contain close pairs]
\label{lem:directions-net-cardinality-bound}
\uses{def:direction-at-end, def:angle-between-directions-theta}

For any $\delta>0$ there is $K=K(\delta)<\infty$ so that for any set
of realized directions at ${\mathcal E}$ of cardinality $K$,
$\ell_1,\ldots,\ell_{K}$, it must be the case that there are $j$ and
$j'$ with $j\not=j'$ such that $\theta(\ell_j,\ell_{j'})<\delta$.
\end{lemma}

\begin{proof}
\sketch
Let $K$ be such that, given $K$ points in the central $2$-sphere of
any $2\epsilon$-tube of scale $1$, at least two are within distance
$\delta/2$ of each other. Now suppose that we have $K$ directions
$\ell_1,\ldots,\ell_K$ at ${\mathcal E}$. Let
$\gamma_1,\ldots,\gamma_K$ be minimizing geodesic rays limiting to
${\mathcal E}$ that represent these directions. Choose a point $x$
sufficiently close to the end ${\mathcal E}$ so that all the
$\gamma_j$ cross the central $2$-sphere $S^2$ of the
$2\epsilon$-neck centered at $x$. By replacing the $\gamma_j$ with
sub-rays we can assume that for each $j$ the endpoint $x_j$ of
$\gamma_j$ lies in $S^2$. Let $d_j$ be the length of $\gamma_j$. By
taking $x$ sufficiently close to ${\mathcal E}$ we can also assume
the following. For each $j$ and $j'$, the angle at $e$ of the
Euclidean triangle $\alpha_je\alpha_{j'}$, where $|\alpha_je|=d_j;
|\alpha_{j'}e|=d_{j'}$ and $|\alpha_j\alpha_{j'}|=d(x_j,x_{j'})$ is
within $\delta/2$ of $\theta(\ell_j,\ell_{j'})$. Now there must be
$j\not= j'$ with $d(x_i,x_j)<(\delta/2)r_i$ where $r_i$ is the scale
of $N_i$. Since $d_j,d_{j'}>\epsilon^{-1}r_i/2$, it follows that the
angle at $e$ of $\alpha_je\alpha_{j'}$ is less than $\delta/2$.
Consequently, $\theta(\ell_j,\ell_{j'})<\delta$.
\end{proof}

Recall that a $\delta$-net in a metric space $X$ is a finite set of
points such that $X$ is contained in the union of the
$\delta$-neighborhoods of these points. The above lemma immediately
yields:

\begin{corollary}[Compactness and nets for the space of directions]
\label{cor:directions-space-compact-net}
\uses{lem:directions-net-cardinality-bound, def:cone-of-directions}
\label{deltanet}

The metric completion $\bar X({\mathcal E})$ of the space of
directions at ${\mathcal E}$ is a compact space. For every
$\delta>0$ this space has a  $\delta$-net consisting of realized
directions. For every $0<r<R<\infty$ the annular region $A_{\mathcal
E}(r,R)=\bar{X}({\mathcal E})\times [r,R]$ in $C_{\mathcal E}$ has a
$\delta$-net consisting of points $(\ell_i,s_i)$ where for each $i$
we have $\ell_i$ is a realizable direction and $r<s_i<R$.
\end{corollary}


\subsection{Completion of the proof of a cone limit at ${\mathcal E}$}

Now we are ready to prove Proposition~\ref{prop:cone-limit-gromov-hausdorff}. In fact, we
prove a version of the proposition that identifies the sequence of
points $x_n$ and also identifies the cone to which the
rescaled manifolds converge.

\begin{proposition}[Explicit Gromov--Hausdorff convergence to the cone]
\label{prop:gromov-hausdorff-limit-is-cone}
\uses{def:cone-of-directions, cor:directions-space-compact-net, prop:incomplete-limit-tube-properties}
\label{secondstatment}

Let $(U_\infty,g_\infty)$  be an incomplete Riemannian $3$-manifold
of non-negative curvature with an end ${\mathcal E}$ as in the
hypothesis of Proposition~\ref{prop:cone-limit-gromov-hausdorff}. Fix a minimizing geodesic
ray $\gamma$ limiting to ${\mathcal E}$. Let $\lambda_n$ be any
sequence of positive numbers tending to infinity. For each $n$
sufficiently large let $x_n\in \gamma$ be the point at distance
$\lambda_n^{-1/2}$ from the end ${\mathcal E}$. Then the based
metric spaces $(U_\infty,\lambda_ng_\infty,x_n)$ converge in the
Gromov-Hausdorff sense to
$\left( C'_{\mathcal E},g_{\mathcal E},([\gamma],1)\right)$. Under
this convergence the distance function from the end ${\mathcal E}$
in $(U_\infty,\lambda_ng_\infty)$ converges to the distance function
from the cone point in the open cone.
\end{proposition}


\begin{proof}
\sketch
\uses{cor:directions-space-compact-net,cor:rescaled-cone-distance-bound,lem:theta-monotone-convergence}
It suffices to prove that given any subsequence of the original
sequence, the result holds for a further subsequence. So let us
replace the given sequence by a subsequence. Recall that for each
$0<r<R<\infty$ we have $A_{\mathcal E}(r,R)\subset C'_{\mathcal E}$,
the compact annulus which is the image of $\bar X({\mathcal
E})\times [r,R]$. The statement about the non-compact spaces
converging in the Gromov-Hausdorff topology, means that for each
compact subspace $K$ of $ C'_{\mathcal E}$ containing the base
point, for all $n$ sufficiently large, there are compact subspaces
$K_n\subset (U_\infty,\lambda_ng_\infty)$ containing $x_n$ with the
property that the $(K_n,x_n)$ converge in the Gromov-Hausdorff
topology to $(K,x)$ (see Section D of Chapter 3, p. 39, of
\cite{Gromov}).

Because of this, it suffices to  fix $0<r<1<R<\infty$ arbitrarily
and prove the convergence result for $A_{\mathcal E}(r,R)$. Since
the Gromov-Hausdorff distance from a compact pointed metric space to
a $\delta$-net in it containing the base point is at most $\delta$,
it suffices to prove that for $\delta>0$ there is a $\delta$-net
$({\mathcal N},h)$ in $A_{\mathcal E}(r,R)$, with $([\gamma],1)\in
{\mathcal N}$ such that for all $n$ sufficiently large there are
embeddings $\varphi_n$ of ${\mathcal N}$ into $A_n(r,R)=\bar
B_{\lambda_ng_\infty}({\mathcal E},R)\setminus
B_{\lambda_ng_\infty}({\mathcal E},r)$  with the following four
properties:
\begin{enumerate}
\item[(1)]  $\varphi_n^*(\lambda_ng_\infty)$ converge to $h$ as $n\rightarrow\infty$,
\item[(2)] $\varphi_n([\gamma],1)=x_n$,
and
\item[(3)]  $\varphi_n({\mathcal N})$ is a $\delta$-net in $A_n(r,R)$, and
\item[(4)] denoting the cone point by $c\in C_{\mathcal E}$,
if $d(p,c)=r$ then $d(\varphi_n(p),{\mathcal E})=r$.
\end{enumerate}

According to Corollary~\ref{cor:directions-space-compact-net} there is a $\delta$-net
${\mathcal N}\subset A_{\mathcal E}(r,R)$ consisting of points
$(\ell_i,s_i)$ where the $\ell_i$ are realizable directions and
$r<s_i< R$. Add $([\gamma],1)$ to ${\mathcal N}$ if necessary so
that we can assume that $([\gamma],1)\in {\mathcal N}$. Let
$\gamma_i$ be a minimizing geodesic realizing $\ell_i$ and let $d_i$
be its length.

Fix $n$ sufficiently large so that $\lambda_n^{-1/2}R\le d_i$ for all $i$. We
define $\varphi_n\colon {\mathcal N}\to A_n(r,R)$ as follows. For any
$a_i=([\gamma_i],s_i)\in {\mathcal N}$ we let
$\varphi_n(a_i)=\gamma_i(\lambda_n^{-1/2}s_i)$. (Since $\lambda_n^{-1/2}s\le
\lambda_n^{-1/2}R\le d_i$, the geodesic $\gamma_i$ is defined at
$\lambda_n^{-1/2}s_i$.) This defines the embeddings $\varphi_n$ for all $n$
sufficiently large. Notice that
$$d_{g_\infty}\left(\varphi_n(\ell_i,s_i),\varphi_n(\ell_j,s_j)\right)=\lambda_n^{-1}s_i^2+\lambda_n^{-1}
s_j^2-2\lambda_n^{-1}s_is_j\theta(\gamma_i,\lambda_n^{-1/2}s_i,\gamma_j,\lambda_n^{-1/2}s_j),$$
or equivalently
$$d_{\lambda_ng_\infty}\left(\varphi_n(\ell_i,s_i),\varphi_n(\ell_j,s_j)\right)=s_i^2+s_j^2-2s_is_j
\theta(\gamma_i,\lambda_n^{-1/2}s_i,\gamma_j,\lambda_n^{-1/2}s_j).$$
Because of the convergence result on angles (Lemma~\ref{lem:theta-monotone-convergence}),
for all $i$ and $j$ we have \begin{eqnarray*} \mathrm{lim}_{n\rightarrow\infty}d_{\lambda_ng_\infty}\left(\varphi_n(\ell_i,s_i),\varphi_n(\ell_j,s_j)\right)
 & = & s_i^2+s_j^2-2s_is_j\mathrm{cos}(\theta(\gamma_i,\gamma_j)) \\ & = & d_{g_{\mathcal
E}}\left((\ell_i,s_i),(\ell_j,s_j)\right).\end{eqnarray*}
 This
establishes the existence of the $\varphi_n$ for all $n$
sufficiently large satisfying the first condition. Clearly, from the
definition $\varphi_n([\gamma],1)=x_n$, and for all $p\in {\mathcal
N}$ we have $d(\varphi_n(p),{\mathcal E})=d(p,c)$.

It remains to check that for all $n$ sufficiently large
$\varphi_n({\mathcal N})$ is a $\delta$-net in $A_n(r,R)$. For $n$
sufficiently large let $z\in A_n(r,R)$ and let $\gamma_z$ be a
minimizing geodesic ray from $z$ to ${\mathcal E}$ parameterized by
the distance from the end. Set
$d_n=d_{\lambda_ng_\infty}(z,{\mathcal E})$, so that $r\le d_n\le
R$. Fix $n$ sufficiently large so that $\lambda_n^{-1/2}R<d_i$ for
all $i$. The point $([\gamma_z],d_n)\in C_{\mathcal E}$ is contained
in $A_{\mathcal E}(r,R)$ and hence there is an element
$a=([\gamma_i],s_i)\in {\mathcal N}$ within distance $\delta$ of
$([\gamma_z],d_n)$ in $C_{\mathcal E}$. Since $s_i\le R$,
$\lambda_n^{-1/2}s_i\le d_i$ and hence
$x=\gamma_i(\lambda_n^{-1/2}s_i)$ is defined. By
Corollary~\ref{cor:rescaled-cone-distance-bound} we have
$$d_{\lambda_ng_\infty}(x,z)\le
d_{g_{\mathcal E}}\left(([\gamma],d_n),([\gamma_i],s_i)\right)\le
\delta.$$ This completes the proof that for $n$ sufficiently large
the image $\varphi_n({\mathcal N})$ is a $\delta$-net in $A_n(r,R)$.

This shows that the $(U_\infty,\lambda_ng_\infty,x_n)$ converge in
the Gromov-Hausdorff topology to $(C'_{\mathcal E},g_{\mathcal
E},([\gamma],1))$.
\end{proof}

\begin{remark}[Non-uniqueness of the incomplete Gromov--Hausdorff limit]
\label{rem:noncomplete-multiple-limits}
Notice that since the manifolds $(U_\infty,\lambda_ng_\infty,x_n)$
are not complete, there can be more than one Gromov-Hausdorff limit.
For example we could take the full cone as a limit. Indeed, the cone
is the only Gromov-Hausdorff limit that is complete as a metric
space.
\end{remark}


\section{Comparison of the Gromov-Hausdorff limit and the smooth
limit}

Let us recap the progress to date. We constructed an incomplete
geometric blow-up limit $({\mathcal U}_\infty,G_\infty,z_\infty)$
for our original sequence. It has non-negative Riemann curvature. We
showed that the zero time-slice $U_\infty$ of the limit is
diffeomorphic to a tube $S^2\times (0,1)$ and that at one end of the
tube the scalar curvature goes to infinity. Also, any point
sufficiently near this end is the center of an evolving
$2\epsilon$-neck defined for an interval of normalized  time of
length $1/2$ in the limiting flow. Then we took a further blow-up
limit. We chose a sequence of points $x_n\in U_\infty$ tending to
the end ${\mathcal E}$ where the scalar curvature goes to infinity.
Then we formed $(U_\infty,\lambda_ng_\infty,x_n)$ where the distance
from $x_n$ to the end ${\mathcal E}$ is $\lambda_n^{-1/2}$. By
fairly general principles (in fact it is a general theorem about
manifolds of non-negative curvature) we showed that this sequence
converges in the Gromov-Hausdorff
sense to a cone.

The next step is to show that this second blow-up limit also exists
as a geometric limit away from the cone point. Take a sequence of
points $x_n\in U_\infty$ tending to ${\mathcal E}$. We let
$\lambda'_n=R(x_n)$, and we consider the based Riemannian manifolds
$(U_\infty,\lambda_n'g_\infty(0),x_n)$. Let $B_n\subset U_n$ be the
metric ball of radius $\epsilon^{-1}/3$ centered at $x_n$ in
$(U_\infty,\lambda'_ng_\infty(0))$. Since this ball is contained in
a $2\epsilon$-neck centered at $x_n$, the curvature on this ball is
bounded, and this ball has compact closure in $U_\infty$. Also, for
each $y\in B_n$, there is a rescaled flow $\lambda'g(t)$  defined on
$\{y\}\times (-1/2,0]$ whose  curvature on $B_n\times (-1/2,0]$ is
bounded. Hence, by Theorem~\ref{sliceconv} we can pass to a
subsequence and extract a geometric limit. In fact, by
Proposition~\ref{partialflowlimit} there is even a geometric
limiting flow defined on the time interval $(-1/2,0]$.

We must compare the zero time-slice of this geometric limiting flow
with the corresponding open subset of the Gromov-Hausdorff
limit constructed in the previous
section. Of course, one obvious difference is that we have used
different blow-up factors: $d(x_n,{\mathcal E})^{-2}$ in the first
case and $R(x_n)$ in the second case. So one important ingredient in
comparing the limits will be to compare these factors, at least in
the limit.


\subsection{Comparison of the blow-up factors}


Now let us compare the two limits: (i) the Gromov-Hausdorff limit of
the sequence $(U_\infty,\lambda_ng_\infty,x_n)$ and (ii) the
geometric limit of the sequence $(U_\infty,\lambda_n'g_\infty,x_n)$
 constructed above.

\begin{claim}[The blow-up ratios are bounded]
\label{claim:blowup-ratio-bounded}
\uses{prop:incomplete-limit-tube-properties, prop:gromov-hausdorff-limit-is-cone, S2intersection, directions}

The ratio $\rho_n=\lambda_n'/\lambda_n$ is bounded above and below
by positive constants.
\end{claim}


\begin{proof}
\sketch
First of all, since there is a $2\epsilon$-neck centered at $x_n$,
by Proposition~\ref{S2intersection} we see that the distance
$\lambda_n^{-1/2}$ from $x_n$ to ${\mathcal E}$ is at least
$R(x_n)^{-1/2}\epsilon^{-1}/2=(\lambda'_n)^{-1/2}\epsilon^{-1}/2$.
Thus,
$$\rho_n^{-1}=\lambda_n/\lambda'_n\le 4\epsilon^2.$$

On the other hand, suppose that
$\rho_n=\lambda_n'/\lambda_n\rightarrow \infty$ as $n\rightarrow
\infty$. Rescale by $\lambda'_n$ so that $R(x_n)=1$. The distance
from $x_n$ to ${\mathcal E}$ is $\sqrt{\rho_n}$. Then by
Lemma~\ref{directions} with respect to this metric there is a sphere
of diameter at most $2\pi$  through  $x_n$ that separates all points
at distance at most $\sqrt{\rho_n}-\epsilon^{-1}$ from ${\mathcal
E}$ from all points at distance at least
$\sqrt{\rho_n}+\epsilon^{-1}$ from ${\mathcal E}$. Now rescale the
metric by $\rho_n$. In the rescaled metric there is a $2$-sphere of
diameter at most $2\pi/\sqrt{\rho_n}$ through $x_n$ that separates
all points at distance at most $1-\epsilon^{-1}/\sqrt{\rho_n}$ from
${\mathcal E}$ from all points at distance at least
$1+\epsilon^{-1}/\sqrt{\rho_n}$ from ${\mathcal E}$. Taking the
Gromov-Hausdorff limit of these spaces, we see that the base point
$x_\infty$ separates all points of distance less than one from
${\mathcal E}$ from all points of distance greater than one from
${\mathcal E}$. This is impossible since the Gromov-Hausdorff limit
is a cone that is not the cone on a single point.
\end{proof}


\subsection{Completion of the comparison of the blow-up limits}

Once we know that the $\lambda_n/\lambda_n'$ are bounded above and
below by positive constants, we can pass to a subsequence so that
these ratios converge to a finite positive limit. This means that
the Gromov-Hausdorff limit of the sequence of based metric spaces
$(U_\infty,\lambda_n'g_\infty,x_n)$ is a  cone, namely the
Gromov-Hausdorff limiting cone constructed is Section~\ref{GromHaus}
rescaled by $\lim_{n\rightarrow\infty}\rho_n$. In particular,
the balls of radius $\epsilon^{-1}/2$ around the base points in this
sequence converge in the Gromov-Hausdorff sense to the ball of
radius $\epsilon^{-1}/2$ about the base point of a cone.


But we have already seen that the balls of radius $\epsilon^{-1}/2$
centered at the base points converge geometrically to a limiting
manifold. That is to say, on every   ball of radius less than
$\epsilon^{-1}/2$ centered at the base point the metrics converge
uniformly in the $C^\infty$-topology to a limiting smooth metric.
Thus, on every ball of radius less than $\epsilon^{-1}/2$ centered
at the base point the limiting smooth metric is isometric to the
metric of the Gromov-Hausdorff limit. This means that the limiting
smooth metric on the ball $B_\infty$ of radius $\epsilon^{-1}/2$
centered at the base point is isometric to an open subset of a cone.
Notice that the scalar curvature of the limiting smooth metric at
the base point is $1$, so that this cone is a non-flat cone.




\section{The final contradiction}

We have now shown that the smooth limit of the balls of radius
$\epsilon^{-1}/2$ centered at the base points of
$(U_\infty,\lambda_n'g_\infty,x_n)$ is isometric to an open subset
of a non-flat cone, and is also the zero time-slice of a
Ricci flow defined for the time interval $(-1/2,0]$. This
contradicts Proposition~\ref{nocones},  one of the consequences of
the maximum principle established by Hamilton. The contradiction
shows that the limit $(U_\infty,g_\infty,x_\infty)$ cannot exist.
The only assumption that we made in order to construct this limit
was that Theorem~\ref{thm:bounded-curvature-bounded-distance} did not hold for some $A_0<\infty$.
Thus, we have established Theorem~\ref{thm:bounded-curvature-bounded-distance}  by contradiction.
